\chapter{Reproducibilidad del canal analítico}
\label{ap:reproducibilidad}

\section{Estructura y orden de ejecución}

El análisis se implementa en seis cuadernos de Jupyter que deben ejecutarse en
orden, más un módulo compartido que fija el contrato entre ellos:

\begin{table}[htbp]
\centering
\caption{Artefactos del canal analítico.}
\label{tab:artefactos}
\small
\setlength{\tabcolsep}{4pt}
\begin{tabular}{p{3.9cm} p{5.4cm} p{4.4cm}}
\toprule
Componente & Función & Salidas \\
\midrule
\texttt{tesis\_common.py} &
Contrato compartido: rutas, cohorte, hiperparámetros, partición, punto de
operación y métricas de equidad & --- \\
\texttt{01\_exploracion\allowbreak\_datos\_dge} &
Ingesta, recodificación, control de fuga, análisis exploratorio &
\texttt{*.parquet}, \texttt{*.manifest.json}, \texttt{calidad\allowbreak\_datos.json},
figura del EDA \\
\texttt{02\_modelo\_referencia} &
Entrenamiento, calibración isotónica, punto de operación, importancia y
sensibilidad al subconjunto hospitalizado y a la exclusión de
\texttt{NEUMONIA} &
\texttt{modelo\_hgb\allowbreak\_mortalidad\allowbreak.joblib},
\texttt{modelo\allowbreak\_referencia.json},
\texttt{sensibilidad\allowbreak\_hospitalizados.json} \\
\texttt{03\_auditoria\_equidad} &
Auditoría por entidad, sexo, edad e interseccional &
\texttt{auditoria\allowbreak\_equidad.json}, figura de equidad \\
\texttt{04\_mitigacion\_pareto} &
Correlatos, tres familias de mitigación, frontera de Pareto &
\texttt{resultados\allowbreak\_tesis.json}, figura de Pareto \\
\texttt{05\_transportabilidad\allowbreak\_loeo} &
Dejar-una-entidad-fuera: \loeoNent{} modelos, cada uno excluyendo una entidad &
\texttt{transportabilidad\allowbreak\_loeo.json}, figura de transportabilidad \\
\texttt{06\_incertidumbre} &
Remuestreo (\boB{} réplicas) y prueba de permutación (\boP) sobre las brechas &
\texttt{incertidumbre\allowbreak\_brechas.json}, figura de incertidumbre \\
\texttt{generar\_tablas.py} &
Traduce los JSON a tablas y macros de \LaTeX{} &
\texttt{tesis/tablas/}, \texttt{macros\allowbreak\_resultados.tex} \\
\bottomrule
\end{tabular}
\end{table}

\section{Contrato de datos entre etapas}

El riesgo característico de un canal por cuadernos es que una etapa consuma un
artefacto producido por una versión anterior de la etapa previa, sin que nada lo
señale. Para evitarlo, cada etapa verifica la procedencia de lo que consume antes
de usarlo:

\begin{enumerate}[leftmargin=*]
  \item El cuaderno 01 escribe, junto al conjunto de datos limpio, un
    \textbf{manifiesto} en JSON con la cohorte empleada, el número de filas, la
    lista de columnas, las tasas de desenlace y las versiones de biblioteca.
  \item Los cuadernos 02 a 04 llaman a \texttt{cargar\_cohorte()}, que exige la
    presencia del manifiesto y verifica columnas requeridas, número de filas y
    bandera de cohorte. Cualquier discrepancia detiene la ejecución con un mensaje
    accionable en lugar de continuar.
  \item El modelo se persiste como \textbf{paquete} --estimador, calibrador,
    lista de predictores, hiperparámetros y huella de la cohorte-- y no como
    estimador desnudo. \texttt{cargar\_modelo()} rechaza un modelo entrenado sobre
    otra cohorte o con otros predictores.
  \item La escritura del conjunto limpio \textbf{falla de inmediato} si no puede
    completarse, en lugar de degradar a otro formato: una degradación silenciosa
    dejaría en disco el artefacto de la corrida anterior.
\end{enumerate}

El cuaderno 05 es la única excepción deliberada a la regla de que solo el 02
entrena: su objeto de estudio \emph{es} un conjunto de modelos que nunca vieron la
entidad evaluada. El 06 no entrena: reconstruye el escenario auditado y solo
remuestrea sobre él. Usa los hiperparámetros canónicos, de modo que la única
diferencia con la línea de base es qué entidades entraron al ajuste.

Los hiperparámetros del modelo, la partición, el punto de operación, los cortes de
edad y las métricas de equidad tienen una única definición, en
\texttt{tesis\_common.py}. Ninguna etapa los redefine, de modo que la alineación
entre cuadernos es estructural y no depende de que coincidan por convención.

\section{Determinismo}

Todas las operaciones estocásticas usan semilla fija (42, salvo la submuestra del
método de in-procesamiento, que usa 3). La partición se deriva de índices
posicionales estratificados por el desenlace, de modo que los bloques de
entrenamiento, calibración y prueba son idénticos en los cuadernos 02, 03 y 04.

Un caso merece mención aparte porque no se resuelve fijando la semilla del ajuste.
La reducción por gradiente exponenciado de Fairlearn devuelve una distribución
sobre clasificadores, y su método de predicción \emph{muestrea} de ella: sin una
semilla en la llamada a predecir, la misma ejecución sobre el mismo modelo ya
ajustado produce brechas distintas. Al detectarlo se midió el efecto --brechas
entre \inMin{} y \inMax{} sobre \inRealizaciones{} realizaciones-- y se fijó la
semilla también en la predicción. Sin esa corrección, el resultado del método de
in-procesamiento no era reproducible entre ejecuciones.

\section{Entorno de ejecución}

\begin{table}[htbp]
\centering
\caption{Versiones utilizadas para producir los resultados de este documento.}
\label{tab:entorno}
\small
\begin{tabular}{l l l l}
\toprule
Componente & Versión & Componente & Versión \\
\midrule
Python        & 3.14.5  & \texttt{scipy}       & 1.18.0 \\
\texttt{numpy}        & 2.5.1   & \texttt{fairlearn}   & 0.14.0 \\
\texttt{pandas}       & 3.0.5   & \texttt{matplotlib}  & 3.11.1 \\
\texttt{scikit-learn} & 1.9.0   & \texttt{joblib}      & 1.5.3 \\
\texttt{pyarrow}      & 25.0.0  & \TeX{} Live          & 2026 \\
\bottomrule
\end{tabular}
\end{table}

\section{Reproducción del documento}

\begin{enumerate}[leftmargin=*]
  \item Colocar el corte anual \texttt{\cohorteFuente} en el directorio raíz, o
    apuntar a él con la variable de entorno \texttt{DGE\_DATA\_PATH}.
  \item Ejecutar los seis cuadernos en orden.
  \item Ejecutar \texttt{python3 generar\_tablas.py} para regenerar tablas y
    macros.
  \item Compilar con \texttt{latexmk -pdf tesis.tex} desde el directorio
    \texttt{tesis/}.
\end{enumerate}

Las cifras del capítulo~\ref{cap:resultados} no se editan a mano en ningún punto:
provienen de macros generadas desde los archivos de resultados. Si una cifra del
texto no coincide con los cuadernos, es porque el paso 3 no se ejecutó, no porque
el texto y el análisis discrepen.
